% GENERATED FILE. Edit canonical lessons, then run npm run generate:books.
% canonical-source-hash: fnv1a64:d355e34de39ca771
% canonical-lessons: MW-C113-lagno, MW-C113-pasand, MW-C113-sauk, MW-C113-het, MW-C113-moh

\chapter{To seem, Liking, Hobby, Affection, Attachment}
\label{ch:mw-to-seem-liking-hobby-affection-attachment}
\hlchaptermodality{113}

\begin{chapteropening}
\textbf{By the end of this chapter:} \emph{I can say to seem, a liking, a hobby, affection and attachment.}

Complete the last lesson of chapter 113: I can say to seem, a liking, a hobby, affection and attachment.
\end{chapteropening}

\section[\texorpdfstring{\mw{लागणो} (lāgṇo)}{लागणो (lāgṇo)}]{\mw{लागणो} (lāgṇo) — to seem, to feel (to someone)}

\label{lesson:MW-C113-lagno}



\begin{quote}
Before the new one: say the Marwadi for an expectation, then the Marwadi for one's will.
\end{quote}

\subsection*{You'll want to know: \mw{लागणो}}
\textbf{\mw{लागणो}} — \emph{lāgṇo} — \textquotedblleft{}to seem, to feel (to someone)\textquotedblright{}.

\begin{cousinweb}
\textbf{\mw{लागणो}} (\emph{lāgṇo}) is a tadbhava, a word inherited from Sanskrit and reshaped by sound change. It comes from \emph{lagyati}, \textquotedblleft{}adheres to\textquotedblright{}. CDIAL lists Old Marwari \emph{lāgaï} under it (CDIAL 10895). The sense moved from \textquotedblleft{}stick to\textquotedblright{} to \textquotedblleft{}seem, feel\textquotedblright{}.
\end{cousinweb}

\begin{grammarlens}[title={What you've built}]
One more word for this chapter.
\end{grammarlens}

\subsection*{Your turn}
\begin{itemize}
\raggedright
  \item \emph{Say it:} \emph{lāgṇo}
  \item \emph{Say it:} \emph{lāgṇo}, once more
  \item \emph{Say it:} say \emph{lāgṇo}
  \item \emph{Recall:} say the Marwadi for an expectation, then the Marwadi for one's will, then say \emph{lāgṇo} again
\end{itemize}

\begin{tcolorbox}[breakable,colback=teal!4,colframe=teal!35!black,title={Before you move on}]
What does \mw{लागणो} mean? (\textquotedblleft{}To seem, to feel (to someone)\textquotedblright{}.) Say it once more.
\end{tcolorbox}
\section[\texorpdfstring{\mw{पसंद} (pasand)}{पसंद (pasand)}]{\mw{पसंद} (pasand) — a liking, a choice}

\label{lesson:MW-C113-pasand}



\begin{quote}
Before the new one: say the Marwadi for one's will, then the Marwadi for to seem.
\end{quote}

\subsection*{You'll want to know: \mw{पसंद}}
\textbf{\mw{पसंद}} — \emph{pasand} — \textquotedblleft{}a liking, a choice\textquotedblright{}.

\begin{cousinweb}
\textbf{\mw{पसंद}} (\emph{pasand}) is a loan, usually traced to Persian \emph{pasand}, \textquotedblleft{}pleasing, approved\textquotedblright{}.
\end{cousinweb}

\begin{grammarlens}[title={What you've built}]
One more word for this chapter.
\end{grammarlens}

\subsection*{Your turn}
\begin{itemize}
\raggedright
  \item \emph{Say it:} \emph{pasand}
  \item \emph{Say it:} \emph{pasand}, once more
  \item \emph{Say it:} say \emph{pasand}
  \item \emph{Recall:} say the Marwadi for one's will, then the Marwadi for to seem, then say \emph{pasand} again
\end{itemize}

\begin{tcolorbox}[breakable,colback=teal!4,colframe=teal!35!black,title={Before you move on}]
What does \mw{पसंद} mean? (\textquotedblleft{}A liking, a choice\textquotedblright{}.) Say it once more.
\end{tcolorbox}
\section[\texorpdfstring{\mw{सौक} (sauk)}{सौक (sauk)}]{\mw{सौक} (sauk) — a hobby, a fondness}

\label{lesson:MW-C113-sauk}



\begin{quote}
Before the new one: say the Marwadi for to seem, then the Marwadi for a liking.
\end{quote}

\subsection*{You'll want to know: \mw{सौक}}
\textbf{\mw{सौक}} — \emph{sauk} — \textquotedblleft{}a hobby, a fondness\textquotedblright{}.

\begin{cousinweb}
\textbf{\mw{सौक}} (\emph{sauk}) is a loan from Arabic \emph{šawq}, \textquotedblleft{}longing, desire\textquotedblright{}, which reached Marwadi through Persian. Its Arabic root is \emph{š-w-q} (Lane). Marwadi says the \emph{š} as \emph{s}.
\end{cousinweb}

\begin{grammarlens}[title={What you've built}]
One more word for this chapter.
\end{grammarlens}

\subsection*{Your turn}
\begin{itemize}
\raggedright
  \item \emph{Say it:} \emph{sauk}
  \item \emph{Say it:} \emph{sauk}, once more
  \item \emph{Say it:} say \emph{sauk}
  \item \emph{Recall:} say the Marwadi for to seem, then the Marwadi for a liking, then say \emph{sauk} again
\end{itemize}

\begin{tcolorbox}[breakable,colback=teal!4,colframe=teal!35!black,title={Before you move on}]
What does \mw{सौक} mean? (\textquotedblleft{}A hobby, a fondness\textquotedblright{}.) Say it once more.
\end{tcolorbox}
\section[\texorpdfstring{\mw{हेत} (het)}{हेत (het)}]{\mw{हेत} (het) — affection}

\label{lesson:MW-C113-het}



\begin{quote}
Before the new one: say the Marwadi for a liking, then the Marwadi for a hobby.
\end{quote}

\subsection*{You'll want to know: \mw{हेत}}
\textbf{\mw{हेत}} — \emph{het} — \textquotedblleft{}affection\textquotedblright{}.

\begin{cousinweb}
The origin of \textbf{\mw{हेत}} (\emph{het}) is not settled. It is often linked with Sanskrit \emph{hita}, \textquotedblleft{}beneficial, kindly\textquotedblright{}, but CDIAL lists \emph{hita} only in the sense \textquotedblleft{}placed\textquotedblright{} (CDIAL 14091) and gives no Marwari or Hindi form like \textbf{\mw{हेत}} (\emph{het}).
\end{cousinweb}

\begin{grammarlens}[title={What you've built}]
One more word for this chapter.
\end{grammarlens}

\subsection*{Your turn}
\begin{itemize}
\raggedright
  \item \emph{Say it:} \emph{het}
  \item \emph{Say it:} \emph{het}, once more
  \item \emph{Say it:} say \emph{het}
  \item \emph{Recall:} say the Marwadi for a liking, then the Marwadi for a hobby, then say \emph{het} again
\end{itemize}

\begin{tcolorbox}[breakable,colback=teal!4,colframe=teal!35!black,title={Before you move on}]
What does \mw{हेत} mean? (\textquotedblleft{}Affection\textquotedblright{}.) Say it once more.
\end{tcolorbox}
\section[\texorpdfstring{\mw{मोह} (moh)}{मोह (moh)}]{\mw{मोह} (moh) — attachment}

\label{lesson:MW-C113-moh}



\begin{quote}
Before the new one: say the Marwadi for a hobby, then the Marwadi for affection.
\end{quote}

\subsection*{You'll want to know: \mw{मोह}}
\textbf{\mw{मोह}} — \emph{moh} — \textquotedblleft{}attachment\textquotedblright{}.

\begin{cousinweb}
\textbf{\mw{मोह}} (\emph{moh}) is a tadbhava, a word inherited from Sanskrit and reshaped by sound change. It comes from \emph{moha}, \textquotedblleft{}bewilderment\textquotedblright{}. CDIAL gives no Marwari form in this entry (CDIAL 10360); the match is through Hindi \emph{moh}. The sense moved from \textquotedblleft{}bewilderment\textquotedblright{} to \textquotedblleft{}infatuation, attachment\textquotedblright{}.
\end{cousinweb}

\begin{grammarlens}[title={What you've built}]
One more word for this chapter.
\end{grammarlens}

\subsection*{Your turn}
\begin{itemize}
\raggedright
  \item \emph{Say it:} \emph{moh}
  \item \emph{Say it:} \emph{moh}, once more
  \item \emph{Say it:} say \emph{moh}
  \item \emph{Recall:} say the Marwadi for a hobby, then the Marwadi for affection, then say \emph{moh} again
\end{itemize}

\begin{tcolorbox}[breakable,colback=teal!4,colframe=teal!35!black,title={Before you move on}]
What does \mw{मोह} mean? (\textquotedblleft{}Attachment\textquotedblright{}.) Say it once more.
\end{tcolorbox}
